\chapter{What is open}\label{chap:open}

\section*{At a glance}
\begin{itemize}
\item \lead{Goal} List what is not done.
\item \lead{Status} \stOpen{} the substitution lemmas \code{Erases.inst} and
  \code{Erases.inst\_let}, \code{KernameInj}, the simulation \code{erases\_correct}, \hl{every
  statement about the traversal and \code{\#erase}}, four divergence entries, and the non-vacuity
  instances other than NV-1's model side and the oracle instances.
\end{itemize}

\section{Planned statements}\label{sec:open-planned}

The planned nodes, per chapter; none of their declarations exists. The table is generated from the
chapters by \code{blueprint/scripts/audit.py}.

% Generated by blueprint/scripts/audit.py --update from the chapters. Do not edit.
No node is planned.


\begin{itemize}
\item \lead{Not stated} the state invariant of \code{erasePure\_erases}, fixed after the review of
  the traversal; the case lemmas of \code{erases\_correct} and \code{erasePure\_erases}; the
  non-vacuity instances other than NV-1's model side (Section~\ref{sec:tests-instances}).
\item \lead{Adversarial reviews} before their proofs: the fixpoint rules of the source semantics
  and \code{ErasesBlock} (DV-7, DV-8); the traversal's frame lemmas and state invariant; the
  closing of recursive blocks.
\item \lead{Unused hypothesis} \code{ErasesDeps.eval} takes \code{hb : BlocksErased}, which its
  proof does not use and MetaRocq's \code{erases\_deps\_eval} does not take; whether the statement
  keeps it is open.
\item \lead{\code{hrun} of NV-1} by symbolic evaluation of the traversal at the pure backend.
\end{itemize}

\section{Roots of the proof package}\label{sec:open-roots}

Every declaration of \code{EraseProof} outside the tests is reachable from a root of
\code{proof/ROOTS.txt}; each root names its planned consumer, the declaration that uses it. The final
form is the single root \code{erase\_correct}. The table is generated from
\code{proof/ROOTS.txt}.

% Generated by blueprint/scripts/audit.py --update from proof/ROOTS.txt. Do not edit.
\begin{tabular}{p{0.4\linewidth}p{0.4\linewidth}p{0.12\linewidth}}
\toprule
Root theorem & Planned consumer & Unit \\
\midrule
\code{EraseProof.\allowbreak{}TrS.\allowbreak{}det} (\ref{lem:EraseProof-TrS-det}) & \code{EraseProof.\allowbreak{}erases\_\allowbreak{}correct} (\ref{thm:EraseProof-erases_correct}) & v4-sim-main \\
\code{EraseProof.\allowbreak{}TrS.\allowbreak{}fvarsIn} (\ref{lem:EraseProof-TrS-fvarsIn}) & \code{EraseProof.\allowbreak{}erasePure\_\allowbreak{}erases} (\ref{thm:EraseProof-erasePure_erases}) & v7-core-main \\
\code{EraseProof.\allowbreak{}TrS.\allowbreak{}uninstantiateN} (\ref{lem:EraseProof-TrS-uninstantiateN}) & \code{EraseProof.\allowbreak{}Erases.\allowbreak{}uninstantiateN} (\ref{lem:EraseProof-Erases-uninstantiateN}) & v3-rel-abstract \\
\code{EraseProof.\allowbreak{}TrS.\allowbreak{}inst\_\allowbreak{}fvar} (\ref{lem:EraseProof-TrS-inst_fvar}) & \code{EraseProof.\allowbreak{}erasePure\_\allowbreak{}erases} (\ref{thm:EraseProof-erasePure_erases}) & v7-core-main \\
\code{EraseProof.\allowbreak{}ProgEnv.\allowbreak{}ordered} (\ref{lem:EraseProof-ProgEnv-ordered}) & \code{EraseProof.\allowbreak{}SrcEval.\allowbreak{}defeq} (\ref{lem:EraseProof-SrcEval-defeq}) & v2-sr-eval \\
\code{EraseProof.\allowbreak{}ProgEnv.\allowbreak{}wf} (\ref{lem:EraseProof-ProgEnv-wf}) & \code{EraseProof.\allowbreak{}SrcEval.\allowbreak{}defeq} (\ref{lem:EraseProof-SrcEval-defeq}) & v2-sr-eval \\
\code{EraseProof.\allowbreak{}ProgEnv.\allowbreak{}lookup} (\ref{lem:EraseProof-ProgEnv-lookup}) & \code{EraseProof.\allowbreak{}ErasableS.\allowbreak{}atom} (\ref{lem:EraseProof-ErasableS-atom}) & v2-erasable-eval \\
\code{EraseProof.\allowbreak{}ProgEnv.\allowbreak{}unfold} (\ref{lem:EraseProof-ProgEnv-unfold}) & \code{EraseProof.\allowbreak{}SrcEval.\allowbreak{}defeq} (\ref{lem:EraseProof-SrcEval-defeq}) & v2-sr-eval \\
\code{EraseProof.\allowbreak{}TrS.\allowbreak{}defeqDFC} (\ref{lem:EraseProof-TrS-defeqDFC}) & \code{EraseProof.\allowbreak{}Erases.\allowbreak{}inst\_\allowbreak{}let} (\ref{lem:EraseProof-Erases-inst_let}) & v3-rel-subst \\
\bottomrule
\end{tabular}



\section{Shipping defects}\label{sec:open-shipping}

\begin{itemize}
\item \lead{Reported, not fixed} the shipping-change register lists the defects it reports without
  fixing them (Section~\ref{reg:reported-not-fixed}).
\item \lead{Peregrine on the pure path} of the 346 \code{validate} and \code{eval} checks on the
  173 in-fragment corpus outputs, 10 fail: unescaped names (R-2), a fixpoint body that is not a
  $\lambda$ (R-11), and an evaluation that reaches an axiom (register entry S-21).
\end{itemize}

\section{Divergence entries}\label{sec:open-divergences}

The register (Chapter~\ref{chap:divergences}) holds the entries of the landed artifacts. The
approved design has these further divergences, each an entry once its artifact lands.

\begin{tabular}{lp{0.3\linewidth}p{0.52\linewidth}}
\toprule
Id & Artifact & What differs \\
\midrule
DV-9 & outputs of \code{erasePure} & not \code{expanded\_eprogram} \\
DV-17 & \code{erasePure}, \code{erase\_correct} & fuel; partial correctness \\
DV-19 & \code{erase\_correct} & no first-order corollary \\
DV-20 & \code{erases\_correct} & Letouzey Theorem 13 big-step; Theorems 12, 15 omitted \\
\bottomrule
\end{tabular}

Register entries whose row in the approved design names artifacts that the entry does not cite:

\begin{tabular}{lp{0.4\linewidth}p{0.4\linewidth}}
\toprule
Id & Landed, not cited & Planned \\
\midrule
DV-3 & \code{SubEnv} & the use of \code{collectDeps} in the core lemmas \\
DV-6 & \code{ProgEnv}, \code{SrcEval}, \code{Erases} & -- \\
DV-11 & -- & the instances NV-5, NV-6, NV-7 \\
DV-13 & -- & \code{closeFix}, \code{closeFix\_substl}, \code{abstract\_eq\_abstract1},
  \code{visit\_agree}, \code{visitAppArgs\_agree}, \code{visitMutual\_agree} \\
DV-14 & -- & \code{KernameInj} \\
\bottomrule
\end{tabular}
